% Claude 1962 proper guide, Fourteenth Sunday after Pentecost: leaf-local
% typesetting primitives, shared by the canonical full research edition and by
% the mechanical synthesis companion compiled from the same source. This file
% owns presentation only. Every textual, historical, exegetical and rights
% judgment lives in the section files and in the leaf's research records.
%
% Web-edition constraint: this leaf declares itself web-eligible, so it uses no
% landscape environment, no multicols and no TikZ. Any ordering a diagram might
% have carried is expressed as an ordered list or a table.

\usepackage{xurl}

\setlength{\emergencystretch}{2.4em}
\widowpenalty=10000
\clubpenalty=10000
\displaywidowpenalty=10000

\newcommand{\work}[1]{\textit{#1}}

% The missal's versicle sign is not in the T1 encoding. It is reproduced as a
% slashed V so the appointed text can be transcribed as printed rather than
% silently normalised to "V.". This formulary prints no response sign.
\usepackage{graphicx}
\newcommand{\Vsign}{\leavevmode\hbox{\rlap{\kern0.16em\raisebox{0.12ex}{\scalebox{0.85}[1.0]{/}}}V}}
\DeclareUnicodeCharacter{2123}{\Vsign}
% The missal prints an acute over the ae-ligature in quǽsumus, Matthǽum and
% quǽrite. T1 carries no precomposed glyph, so the accent is composed over the
% ligature rather than the word being silently respelled.
\DeclareUnicodeCharacter{01FD}{\'{\ae}}
\DeclareUnicodeCharacter{01FC}{\'{\AE}}

% The Gospel heading's Maltese cross, likewise outside T1.
\usepackage{amsmath,amsbsy}
\newcommand{\gospelcross}{\ensuremath{\boldsymbol{+}}}
\DeclareUnicodeCharacter{2720}{\gospelcross}

% Quiet parenthetical proper cues, in liturgical order, for later headings and
% cross-references. Defined once so the abbreviations stay uniform.
\newcommand{\cue}[1]{\textnormal{(\textit{#1})}}

% Keeps an element's heading, identification line and text block together.
\newcommand{\elementguard}{\Needspace{12\baselineskip}}

% Element identification line: Latin incipit, printed heading and reference,
% and the missal's marginal number.
\newcommand{\elementid}[3]{%
  \par{\small\latin{#1} \textperiodcentered{} #2 \textperiodcentered{} marginal no.~#3}\par}

% The complete appointed Latin, inside the full edition only.
\newenvironment{appointedlatin}[1]{%
  \begin{massbox}{#1}\itshape}{\end{massbox}}

% An identified public-domain English witness, never a project translation.
% The title always names the witness and the exact locus it was taken from.
\newenvironment{namedtranslation}[1]{%
  \begin{sourcecard}{#1}\small}{\end{sourcecard}}

% A place where the registered public-domain English does not answer the
% appointed Latin. The gap is stated, never filled.
\newcommand{\englishgap}[1]{\par\textbf{Where the public-domain English does
not answer the Latin.} #1\par}

% A witness dossier inside the detailed commentary: who, at what exact locus,
% arguing what, and in what evidence state. The left rule keeps it distinct
% from the surrounding synthesis without implying doctrinal weight.
\newenvironment{witnessnote}{\begin{studyframe}\small}{\end{studyframe}}
\newcommand{\wlocus}[1]{\textbf{#1}\quad}

% A settled point of apparatus the reader needs in order to use the citations
% above: a numbering convention, a rubrical locus. The title names the subject,
% not the guide's handling of it.
\newcommand{\correction}[2]{%
  \begin{studybox}{#1}#2\end{studybox}}

% The exploratory-proposal notice and the proposals themselves. Every proposal
% renders its anchors, mechanism, fruit and controlling limit as labelled
% fields so no reader can mistake a proposal for a sourced finding.
\newenvironment{proposal}[1]{%
  \begin{massbox}{Editorial proposal --- #1}\small}{\end{massbox}}
\newcommand{\pfield}[2]{\par\textbf{#1:} #2}

% Notable-and-quotable gallery entry: appointed phrase, later user, locus and
% the turn in register. Breakable, so a long entry splits rather than jumping
% whole to the next page and leaving a sparse one behind.
\newenvironment{afterlife}[1]{%
  \begin{massbox}{#1}\small}{\end{massbox}}

% Two-tier page-2 dossier table. This sheet makes three departures from the
% baseline page-2 measurements it was built on, and the note below
% \dossierprose states them together with the escape clause that licenses
% them. The baseline is historical. The 1962 profile stated these figures
% until 2026-09-23 (a9086b3da) and now states none: it names the dossiertable
% and concisedossiertable environments of src/common/propers-format.tex as
% the owner of the measurements. This schema-1 leaf does not import that file
% and keeps its own table here. The first departure is the summary-tier column
% widths set here. The baseline set them at
% 0.14 / 0.18 / 0.35 / 0.16\linewidth; this table sets
% 0.11 / 0.15 / 0.19 / 0.38\linewidth, the four still totalling 0.83 as the
% baseline's did. The movement is into Date, because the corpus answers the
% Gospel and the Communion with six composition claims apiece and the Epistle
% with five, the Offertory and Alleluia also carry a traditional-attribution
% group, and the generated annotation prints them in one cell. Location
% gives up the most, its detail belonging to the explanatory row beneath. The
% summary tier keeps \footnotesize, \arraystretch, \LTpre and \LTpost exactly
% as the baseline set them. An earlier form of this sheet set
% 0.13 / 0.16 / 0.22 / 0.32, when the Date cell printed every source label in
% full, and then 0.13 / 0.17 / 0.24 / 0.29 once the projection printed the
% corpus's concise display labels. The traditional-attribution relation
% lengthened three Date cells again, and the present widths answer that.
\newenvironment{dossiertable}
{\footnotesize
 \renewcommand{\arraystretch}{1.02}%
 \setlength{\LTpre}{0.15em}\setlength{\LTpost}{0pt}%
 \begin{longtable}{@{}>{\raggedright\arraybackslash}p{0.11\linewidth}>{\raggedright\arraybackslash}p{0.15\linewidth}>{\raggedright\arraybackslash}p{0.19\linewidth}>{\raggedright\arraybackslash}p{0.38\linewidth}@{}}
 \toprule
 \textbf{Proper} & \textbf{Citation} & \textbf{Location} & \textbf{Date}\\
 \midrule
 \endhead}
{\bottomrule\end{longtable}}

% The page-2 Date column. Every biblical date printed in this guide comes from
% the Scripture chronology corpus, carried in research/chronology.toml and
% projected into research/chronology-annotations.tex, which owns
% \chronologyannotation, \chronologyannotationgroup and
% \chronologyannotationclaim. This file defines the Date-cell wrapper and
% nothing else, so no date reaches the page except through the generated
% projection.
%
% \chronodate{element-key}{content} -- one dossier row's Date cell.
% Typesets the content; the key renders nothing.
\newcommand{\chronodate}[2]{#2}

% Full-width explanatory and event rows inside a dossier.
%
% Measurement note. THREE of the baseline page-2 measurements (historical; see
% the note above \dossiertable) are varied on this sheet, and these are all of
% them:
%
%   D1. Summary-tier column widths, set above at \dossiertable. Baseline
%       0.14 / 0.18 / 0.35 / 0.16\linewidth; here
%       0.11 / 0.15 / 0.19 / 0.38\linewidth.
%   D2. Spanning-row font size. The baseline set the table in \footnotesize
%       and specified no separate size for the spanning row, so its value
%       for that row is \footnotesize --- 9pt on 11pt leading under the 11pt
%       base class in src/common/preamble.tex. Here it is \fontsize{6.6}{7.2}:
%       a reduction of about 27 percent in glyph size and 35 percent in
%       leading, a little below \scriptsize. This is a large departure and is
%       recorded as one.
%   D3. Spanning-row width. Baseline 0.92\linewidth; here 0.93\linewidth.
%
% The profile provides for all three and still states the condition: change
% measurements only when the complete inventory cannot remain legible after the
% hierarchy and the evidence have been preserved, and never spill, omit
% evidence, or move apparatus here from elsewhere. That is this sheet's case.
% Seven dossiers must stand on one physical page; every explanatory row carries
% a named bounded negative -- a psalm with no per-psalm modern critical dating,
% a witness whose rights are unresolved, an attribution that is the
% commentator's own reconstruction -- and a longtable row cannot break, so the
% sheet either fits whole or starts a second page the profile forbids.
%
% An earlier revision set the spanning rows at 0.98\linewidth, the widths at
% 0.13 / 0.16 / 0.22 / 0.32 and the type size at \fontsize{6.4}{6.9}; the
% chronology projection's concise display labels let all three relax, and the
% spanning rows went back to 0.92. The traditional-attribution relation of
% 2026-09-09 then lengthened the Offertory, Alleluia and Gospel Date cells and
% the Offertory explanation, and the Epistle's dossier no longer fitted on
% page 2. The widths of D1 alone cannot absorb it, since Citation and Location
% then set the height of their rows, so D3 returns at 0.93, a hundredth of the
% line and far short of the earlier 0.98.
%
% What the profile's hierarchy and the baseline fix, and this sheet leaves
% alone: the two-tier hierarchy,
% the four summary fields, the \footnotesize summary tier, \arraystretch 1.02,
% \LTpre 0.15em, \LTpost 0pt, the 0.1em breathing space ending ordinary dossier
% rows, \cmidrule(lr){1-4} within a dossier with \midrule only between complete
% dossiers, the italic event row set between two \cmidrule(lr){1-4} rules ahead
% of its explanation.
\newcommand{\dossierrowsize}{\fontsize{6.6}{7.2}\selectfont}
\newcommand{\dossierprose}[1]{\multicolumn{4}{@{}>{\raggedright\arraybackslash\dossierrowsize}p{0.93\linewidth}@{}}{#1}\\[0.1em]}
\newcommand{\dossierevent}[1]{\multicolumn{4}{@{}>{\raggedright\arraybackslash\dossierrowsize}p{0.93\linewidth}@{}}{\textit{#1}}\\}

% Divergence table used in the commentary: appointed form against the Bible
% text the guide's English is quoted from.
\newenvironment{divergencetable}
{\footnotesize\begin{longtable}{>{\raggedright\arraybackslash}p{0.17\linewidth}>{\raggedright\arraybackslash}p{0.36\linewidth}>{\raggedright\arraybackslash}p{0.38\linewidth}}
\toprule
\textbf{Proper and locus} & \textbf{Appointed} & \textbf{Clementine}\\
\midrule
\endhead}
{\bottomrule\end{longtable}}

% A two-column comparison, for a matrix whose second column carries the whole
% argument. The three-column form with an empty header would leave a third of
% the measure blank and unlabelled.
\newenvironment{pairtable}[2]
{\footnotesize\begin{longtable}{>{\raggedright\arraybackslash}p{0.20\linewidth}>{\raggedright\arraybackslash}p{0.74\linewidth}}
\toprule
\textbf{#1} & \textbf{#2}\\
\midrule
\endhead}
{\bottomrule\end{longtable}}

% Reception and comparison tables used in the detailed commentary.
\newenvironment{receptiontable}[3]
{\footnotesize\begin{longtable}{>{\raggedright\arraybackslash}p{0.20\linewidth}>{\raggedright\arraybackslash}p{0.34\linewidth}>{\raggedright\arraybackslash}p{0.38\linewidth}}
\toprule
\textbf{#1} & \textbf{#2} & \textbf{#3}\\
\midrule
\endhead}
{\bottomrule\end{longtable}}

% A four-column measurement table, for the oration word-counts, whose columns
% are all short.
\newenvironment{counttable}[3]
{\footnotesize\begin{longtable}{>{\raggedright\arraybackslash}p{0.26\linewidth}>{\raggedright\arraybackslash}p{0.14\linewidth}>{\raggedright\arraybackslash}p{0.52\linewidth}}
\toprule
\textbf{#1} & \textbf{#2} & \textbf{#3}\\
\midrule
\endhead}
{\bottomrule\end{longtable}}

% The two outputs share every section file except the appointed text and the
% element-by-element sweep, and no macro here states that difference to a
% reader. What each edition carries and omits is said once, in the terminal
% Appendix: Scope and Qualifications, which both editions print. An
% \editionnote macro was defined in both arms of this branch and invoked
% nowhere; it is withdrawn rather than left as a second, unpublished statement
% of the same fact.

% \sectionguard and the deepstudy reading size come from the shared preamble
% and are deliberately not redefined here.
